\subsection{Incident energy needed for a bounded receiver set}
The source test has a useful design consequence when coherent mixtures of the dictionary are available. Assume $K_{\rm mat}$ invertible, set $q=\nn(K_{\rm mat})$, $Q_s=C_{\mathcal B}-N^*K_{\rm mat}^{-1}N$ and let the incident-energy Gram $W_{\rm in}$ be positive definite on the allowed dictionary. Energy is $\tr(V^*W_{\rm in}V)$. We keep the old error majorant $\Omega=\sigma^2 I_s$, $\sigma>0$, fixed while choosing $V$. Write $\beta_1\ge\cdots\ge\beta_m$ for the eigenvalues of $W_{\rm in}^{-1/2}Q_sW_{\rm in}^{-1/2}$.

\begin{corollary}[Energy needed to remove the source deficit]\label{cor:energy}
Suppose that one material in the class has a physical solution on the full incident dictionary and $q>0$. If $m<q$, $s<q$, or the available $q$th eigenvalue satisfies $\beta_q\le0$, no energy budget removes $h_{\rm src}$. Otherwise the infimum of the energy needed for $h_{\rm src}=0$ is
\begin{equation}
 E_{\rm inf}=\sigma^2\sum_{i=1}^{q}\beta_i^{-1}.
 \label{eq:energy}
\end{equation}
Every budget strictly greater than this infimum suffices. At $q=0$ the source deficit is always zero.
\end{corollary}
\begin{proof}
Applying the material graph argument with $V=I$ and zero error gives $\np(Q_s)\le q$. By \eqref{eq:indexsmall}, zero deficit requires a $q$-dimensional beam subspace on which $V^*Q_sV\succ\sigma^2 I$. Restrict to orthonormal source coordinates on that subspace and write the whitened restricted synthesis matrix as $UT$, with $U^*U=I_q$ and $T$ invertible. Put $A=U^*W_{\rm in}^{-1/2}Q_sW_{\rm in}^{-1/2}U$. Then $T^*AT\succ\sigma^2 I$ implies $TT^*\succ\sigma^2 A^{-1}$. The energy exceeds $\sigma^2\tr A^{-1}\ge E_{\rm inf}$ by eigenvalue interlacing. Choosing the leading generalized eigenbeams with energies just above $\sigma^2/\beta_i$ attains the infimum in the limit.
\end{proof}
This is a requirement for the fixed quadratic prediction to bound every receiving direction at saturation. It does not guarantee saturation, separation of two material classes, or a particular material reconstruction. A finite menu of physical angles can require more energy. An exposure-dependent error model requires a new calculation.
